Combine \wikilink{s23--b001}{the free-product universal property} with
\wikilink{s13--b006}{abelianization} and
\wikilink{s17--b010}{the direct-sum universal property}.

\begin{theorem}\label{free-product-abelianization}
There is a natural isomorphism
$(G*H)^{\mathrm{ab}}\cong G^{\mathrm{ab}}\oplus H^{\mathrm{ab}}$.
\end{theorem}
\begin{proof}\label{free-product-abelianization-proof}
The maps from $G$ and $H$ to the two summands induce a map from the
free product to their abelian direct sum, hence a map
$u:(G*H)^{\mathrm{ab}}\to G^{\mathrm{ab}}\oplus H^{\mathrm{ab}}$.
In the opposite direction, the two maps into $(G*H)^{\mathrm{ab}}$
factor through $G^{\mathrm{ab}}$ and $H^{\mathrm{ab}}$ and then induce
$v$ from their direct sum. The composite $uv$ is the identity on
each summand, hence on the direct sum. The composite $vu$ agrees
with the identity after composition with the maps from $G$ and $H$.
Free-product uniqueness and the surjectivity of abelianization make
$vu$ the identity too. The constructions use the canonical maps only,
which also proves naturality under homomorphisms of the two factors.
\end{proof}
\begin{remark}\label{abelian-pushout-link}
The same method constructs pushouts in abelian groups by taking a
direct sum and imposing the two common-source images as relations.
Their module-theoretic treatment belongs with direct sums and quotients
in \wikilink{a2s09}{the module notes}; the free product itself is the
coproduct in groups and is generally not abelian.
\end{remark}
